\chapter{Conclusão}
\label{cap:conclusao}

% figuras estão no subdiretório "figuras/" dentro deste capítulo
\graphicspath{\currfiledir/figuras/}


Neste trabalho explicamos a intuição por trás da síntese de imagens realista e
como podemos capturar essa intuição com a Equação da Renderização. Em seguida,
descrevemos como podemos usar o método de Monte Carlo para a aproximar a
solução da Equação, usando a técnica de Amostragem de Importância. Descrevemos
duas técnicas de amostragem e como combinar elas utillizando a técnica de
Amostragem de Importância Múltipla. Por fim explicamos como podemos formular a
Equação da Renderização em uma integral de caminhos, possibilitando a
amostragem de caminhos inteiros.

Baseando-nos em um Ray Tracer pronto, implementamos um Path Tracer que utiliza
das técnicas descritas anteriormente, formulando o estimador de Monte Carlo e
os pesos da Amostragem de Importância Múltipla. Finalmente, analisamos os
resultados do Path Tracer, considerando as técnicas, as imagens de variância
geradas, e cenas com diferentes iluminações.

Os resultados obtidos mostraram a complementariedade das técnicas de amostragem
apresentadas, com a Amostragem por BRDF mais eficiente em materiais reflexivos
e fontes de luz grandes, enquanto a Amostragem de Fontes de Luz apresentou
melhores resultados em materiais foscos e fontes de luz pequenas. Combinando as
duas técnicas, obtivemos uma imagem com significativamente menos ruído que as
outras técnicas separadamente, ilustrando as vantagens previstas pela teoria.

Com base na implementação, existem diversas possibilidades de extensão deste
trabalho. \citet{lindqvist21} mostra, entre outras, uma técnica de amostragem
simples baseada apenas no hemisfério ponderado pelo cosseno do ângulo de
incidência, sem considerar o material. Essa estratégia poderia, possivelmente,
ajudar a diminuir a variância da esfera vermelha observada no
\autoref{cap:resultados}.

Dada a formulação com base na Amostragem de Importância Múltipla, também é
possível implementar algoritmos mais complexos como o Path Tracing
Bidirecional, \emph{Metropolis Light Transport} \citep{veach97}, \emph{Vertex
Connection and Merging} \citep{georgiev12}, e mais recentemente, ReSTIR
(\emph{Spatiotemporal Reservoir Importance Resampling}) \citep{bitterli20};
todos esses que usam a Amostragem de Importância Múltipla.

Além disso, ortogonalmente a esses algoritmos, é possível mover a implementação
para a GPU, utilizando \emph{frameworks} de computação paralela, o que
possibilita transformar a implementação em um \emph{Path Tracer} de tempo real.


